\scanpage{64}
Supposons donc les $F_{X'}$ qu. coh., \deleted{\unsure{et}}\inserted{et s'induisant bien si $X''\subset X'$.} De plus, supposons d'abord $X$ affine, donc les $Y(i)$ affines. \deleted{\unsure{les $H$}} Comme on a encore
\[
H^*(\widetilde{X'},F)=H^*(X'_{\mathrm{zar}},F_{X'})
\]
on trouve encore $E_2^{pq}=0$ si $\unsure{q}\ne0$ d'où
\[
H^*(X_{\mathrm{cris}},F)\simeq H^*(\nu\longmapsto F(\underbrace{Y^{\nu+1}_{/X}}_{\substack{\text{complété formel,}\\\text{le long de }X}}))
\]
\margin{\deleted{[\illegible{1} \unsure{de $X$} \unsure{$(i)$} \illegible{2} \unsure{!}]}}
Lorsque $X$ n'est pas affine, on introduit encore les faisceaux
\[
\mathcal F^{(\nu)}:\quad U\longmapsto\varprojlim F(U,U(i)^{\nu+1})
\]
sur $X_{\mathrm{zar}}$ [$U(i)$ est le $i$ième voisinage inf de $U$ dans $Y^{\nu+1}$ ; si $V$ est un ouvert de $Y$ induisant $U$, cela revient à prendre
\[
\mathcal F^{(\nu)}(U)=F(\underbrace{\widehat{V^{\nu+1}}_{/U}}_{\substack{\text{complété formel}\\\text{le long de }U}})
\]
];
Pour $\nu$ variable, on trouve un faisceau cosimplicial, \deleted{\unsure{d'un}} noté $\mathcal F^*$, et on trouve
\mspage{63}
